\chapter{The Two Outputs}
\label{chap:the-two-outputs}

The descent has produced two things: a tape and a particle. This chapter says what they are together, and
the saying is the thesis the whole book was built to reach. Along the way it does one more piece of proved
work --- it places two electrons side by side and reads what their pairing is --- and then it closes the
movement on the sentence the volume exists to state. What comes before this chapter is built; what this
chapter adds is the reading that ties the movement's two products into one claim, offered as a reading and
resting on a great deal that is proved.

\section{One program, three carriers}
\label{se:one-program-three-carriers}

There is one program. The thirty-five-rung tower the bootstrap built is a single construction, and what it
computes is set entirely by the carrier it is run over. Over its own values, it builds itself: the apparatus
of the first movement, each rung constructing the next. Over its own truth, it reads itself and collapses:
true and false made one, the apex of the second movement. Walked back down over number --- traced,
executed, counted --- it becomes the particle, the number in its bracket with its four faces. One program,
three carriers, three outputs, and the construction \emph{tanges} its output by the carrier: the same
machine throughout, the carrier alone deciding what comes out. This is the spine of the book and the idea the
first book does not have. The first book built an instrument and pointed it at the world; this one builds a
single self-hosting program and changes only what it is run over, and watches the output change from a
machine to a collapse to a particle. Carrier-polymorphism --- one program, output set by the carrier --- is
the whole of it, and the loop is closed: what went out over values and collapsed over truth comes back, over
number, as matter.

\section{The electron algebra and the Cooper pair}
\label{se:electron-algebra-cooper-pair}

Because the descent read the particle as a number, two of them can be placed side by side --- and a
class-instance cannot. Ask the tower for its electron twice and you get the one instance back, with nothing
beside it to compare; but two numbers can be put in a product, and ``what do two electrons do together?''
becomes a question the moment the electron is a number. The answer is a commutator: two electron-numbers
combined either commute or they do not. When they commute they \emph{funge}, pooling into matter that stays
matter with no residue; when they do not they \emph{tange} out a residue, a single mark of a flip, the turn
toward the mirror. Here, unlike most of the descent, the meter that counts those flips is proved. Read the
gauge as a corner of a hypercube, each bit a commute or a crank, and the meter is the Hamming distance ---
the number of bits at which two readings differ. On this meter the results are decided: the electron sits at
distance zero, the all-funge corner, the unit that reads as itself and carries no flip; and the distance
counts the cranks, the positrons a reading carries off that unit.

Now bind two of them. Take one electron; take the commutator of a second with it; commute that with the
first again --- a nested product, an electron multiplied into the product of another with itself. The
construction names this the pair, and the proved meter places it: the bound reading sits at distance one, and
that is decided, not designed. A distance-one reading carries exactly one crank, one positron's worth of the
flip, so the pair is matter carrying a single turn toward the mirror --- not the pure matter of the electron
at distance zero, and not a free positron, but two electrons locked into one object that carries, between
them, a single unit of the flip. This is the shape of a superconducting pair, and the construction reaches
the shape and fences it as one: the form of two electrons bound at distance one, with no claim on the theory
the name carries. Two things are kept apart here, as they have been all the way down. The word
\emph{electron} names the unit of this algebra --- the distance-zero corner, proved --- and that use is
earned. What it does not name, here, is the type-level identity that would fix this number as that particle
and no other; that identification is a further claim, reached for and named only at the close. The algebra
names a unit; the naming of the identity waits for the ledger.

\section{The two outputs}
\label{se:the-two-outputs}

On its way down the descent emitted two things, and they are the two the whole construction was built to
produce. The first is a tape: the program the machine runs on itself, which is code. The second is a
particle: the electron, and the pair, which is matter. The decisive thing --- the thing the entire movement
was arranged to make visible --- is that these are not two separate constructions. They are two readings of
one descent: the same backward pass that wrote the tape read off the electron. Code and matter are
\emph{funged} here into one process, the two outputs of a single self-compilation.

Said plainly, and this is the sentence the volume exists to reach: matter and code are the two outputs of one
self-compilation. Matter is what the self-compiler computes when it is run over number; code is what it
computes when it compiles itself; and they are the same machine's output over two carriers, not two unrelated
facts that happen to share a vocabulary. A construction that builds itself, points at its own truth,
collapses, and walks back down emits a program and a particle by the one act of descending --- so the
deepest thing the book has to say about matter is that it is what code looks like when a self-compiler is run
over number, and the deepest thing it has to say about code is that it is what matter looks like when the
same compiler is run over itself.

\section{What is demonstrated, and what is assumed}
\label{se:demonstrated-assumed-III7}

What is demonstrated in this chapter is the electron algebra's proved meter --- the Hamming distance on the
gauge read as a hypercube, the electron at distance zero, the bound pair at distance one, the distance
counting the cranks --- and, drawing on the whole movement, the two outputs the descent produced. The thesis
itself, that matter and code are the two outputs of one self-compilation, is a reading: the synthesis the
construction was arranged to make available, resting on a descent that is built --- the backward pass, the
seam, the trace, the slip test, the three trips, the four faces, all proved --- and offered as exactly what
it is, a reading earned from a great deal that is proved and not a leap made over an empty gap.

What is assumed is unchanged. The seed and its decidability; the honest import, marked as borrowed; the
single standing grant, the law of motion the first book named at its summit, still out of play; and the
theory of superconductivity, named, fenced as a shape. One naming, and only one, is still held back: the
type-level identity that would fix the number the descent recovered as the electron and no other. Everything
that carries it there is built; the naming of the invariant it lands on is the one thing reached for and not
yet grasped, and the movement that follows will grasp it, or say exactly why it is left. With the two outputs
named and the thesis stated, the return is complete: the machine has built itself, read itself, collapsed,
and walked back down --- traced, executed, and counted --- to emit a tape and a particle. What remains is to
say, of the whole construction, what stands proved and what is still reached for, and to set out the laws the
building taught. Those are the last two movements.
